%% paper2.tex
%% IEEE Journal Paper — Research and Evaluation
%% A Comparative Study of RSA Optimization Techniques for Secure Web Applications

\documentclass[journal]{IEEEtran}

\usepackage{cite}
\usepackage{amsmath,amssymb,amsfonts}
\usepackage{graphicx}
\usepackage{textcomp}
\usepackage{booktabs}
\usepackage{multirow}
\usepackage{array}
\usepackage{url}
\usepackage{hyperref}
\usepackage{longtable}
\usepackage{listings}
\usepackage{xcolor}
\usepackage{float}
\usepackage{caption}
\lstdefinestyle{pythonstyle}{
    language=Python,
    basicstyle=\ttfamily\tiny,
    keywordstyle=\bfseries\color{blue!70!black},
    commentstyle=\itshape\color{green!50!black},
    stringstyle=\color{red!60!black},
    numbers=left,
    numberstyle=\tiny\color{gray},
    numbersep=5pt,
    frame=single,
    breaklines=true,
    showstringspaces=false,
    tabsize=4,
    columns=fullflexible,
    float=H,
}
\lstset{style=pythonstyle}
\hypersetup{colorlinks=true,linkcolor=blue,citecolor=blue,urlcolor=blue}

\hyphenation{op-tical net-works semi-conduc-tor}

\begin{document}

\title{A Comparative Study of RSA Optimization Techniques for Secure Web Applications}

\author{
\begin{tabular}{ccc}
Patil~Bhavesh~S. & Patil~Devendra~V. & Vala Shakti P. \\
240093116002 & 230090116054 & 230090116065 \\
CKPCET, Surat & CKPCET, Surat & CKPCET, Surat \\
\end{tabular}
}

\markboth{IEEE Research and Evaluation: RSA Optimization for Web Security}%
{Patil \MakeLowercase{\textit{et al.}}: RSA Optimization Comparative Study}

\IEEEspecialpapernotice{(Research and Evaluation)}

\maketitle

\begin{abstract}
The Rivest--Shamir--Adleman (RSA) cryptosystem remains a cornerstone of public-key infrastructure, yet its computational cost poses significant challenges for modern web applications requiring low-latency secure communication. This paper presents a comparative study of established RSA optimization techniques drawn from over four decades of cryptographic research, evaluating their practical impact on web application performance. We survey and compare algorithms including the Chinese Remainder Theorem (CRT) for accelerated decryption, Montgomery multiplication for efficient modular arithmetic, sliding-window and $k$-ary modular exponentiation, key caching strategies, batch processing, and hybrid RSA+AES envelope encryption. An experimental framework implementing both a standard (textbook) RSA baseline and an optimized variant incorporating all surveyed techniques is evaluated across six benchmark dimensions: correctness, speed, memory consumption, key-size scaling (1024--4096~bits), hybrid versus pure RSA encryption (1~KB--1~MB), and concurrent throughput (1--50 threads). Fresh experimental results demonstrate that key caching yields the greatest practical speedup ($1112\times$ for encryption, $123\times$ for decryption), hybrid RSA+AES-GCM outperforms pure RSA by $5\times$ to $3903\times$ for bulk data, and CRT-based decryption provides measurable gains at 2048-bit keys ($1.27\times$). The comparative analysis synthesizes findings from 30~key references, identifies optimization trade-offs specific to interpreted language environments, and provides actionable recommendations for web application architects.
\end{abstract}

\begin{IEEEkeywords}
RSA Algorithm, Web Application Security, Performance Optimization, Chinese Remainder Theorem, Montgomery Multiplication, Hybrid Encryption, Comparative Study, Post-Quantum Cryptography.
\end{IEEEkeywords}

\IEEEpeerreviewmaketitle


\section{Introduction}
\label{sec:intro}
\IEEEPARstart{T}{he} proliferation of web-based services---spanning electronic commerce, online banking, social media, healthcare information systems, and government portals---has made robust cryptographic mechanisms an essential requirement for protecting data in transit and at rest. Among the many cryptographic algorithms available, the Rivest--Shamir--Adleman (RSA) cryptosystem~\cite{rivest1978} remains one of the most widely deployed asymmetric encryption schemes, underpinning digital signatures, key exchange protocols (TLS/SSL), certificate authorities, and authentication systems worldwide.

RSA's security rests on the computational hardness of the integer factorization problem: given a modulus $n = p \cdot q$ where $p$ and $q$ are large primes, no efficient classical algorithm is known for recovering $p$ and $q$ from $n$ alone~\cite{menezes2001}. However, this mathematical elegance comes at a significant computational cost. RSA operations---particularly decryption and key generation---require modular exponentiation with operands spanning thousands of bits, making them orders of magnitude slower than symmetric alternatives such as AES~\cite{nist_fips197}.

The performance challenge is compounded by the evolving threat landscape. Classical cryptanalytic attacks~\cite{boneh1999}, lattice-based methods~\cite{coppersmith1997}, timing attacks~\cite{kocher1996}, fault injection attacks~\cite{boneh2001}, and the looming threat of quantum computing via Shor's algorithm~\cite{shor1997} all constrain RSA's practical deployment. While post-quantum alternatives such as CRYSTALS-Kyber~\cite{kyber2023} and CRYSTALS-Dilithium~\cite{dilithium2022} are being standardized by NIST~\cite{nist_pqc}, RSA will remain in service for years due to legacy infrastructure, certificate ecosystems, and backward compatibility requirements~\cite{nist_sp80057}.

This paper addresses the practical question: \textit{which RSA optimization techniques, drawn from the existing literature, provide the most significant performance improvements in the context of web applications?} We conduct a systematic comparative study of the major optimization approaches, implement them in a controlled experimental framework, and measure their individual and combined effects using fresh benchmarks.


\section{Literature Survey}
\label{sec:litreview}

This section surveys the key contributions to RSA optimization, security analysis, and hybrid encryption architectures that inform our comparative study.

\subsection{Foundations of Public-Key Cryptography}

The conceptual foundation for RSA was laid by Diffie and Hellman~\cite{diffie1976}, who proposed the concept of public-key cryptography in 1976. Shortly thereafter, Rivest, Shamir, and Adleman~\cite{rivest1978} published the first practical public-key encryption and digital signature scheme, based on the difficulty of factoring the product of two large prime numbers. The RSA algorithm computes a public key $(e, n)$ and private key $(d, n)$ such that encryption is $c = m^e \bmod n$ and decryption is $m = c^d \bmod n$. The security of this scheme was analyzed by Rabin~\cite{rabin1979}, who showed that recovering the plaintext from the ciphertext is as hard as factoring.

Alternative public-key systems have since been proposed, including the ElGamal cryptosystem~\cite{elgamal1985} based on discrete logarithms, and elliptic curve cryptography (ECC) independently proposed by Koblitz~\cite{koblitz1987} and Miller~\cite{miller1985}. While ECC offers equivalent security with smaller key sizes, RSA remains dominant in PKI infrastructure and TLS deployments.

\subsection{RSA Performance Optimization Techniques}

\textbf{Chinese Remainder Theorem (CRT).}
Garner~\cite{garner1959} introduced the residue number system and recombination algorithm that forms the basis of CRT-accelerated RSA decryption. Quisquater and Couvreur~\cite{quisquater1982} were the first to apply CRT to RSA, showing that decryption can be decomposed into two half-size modular exponentiations, yielding a theoretical $4\times$ speedup. The method pre-computes $d_p = d \bmod (p-1)$ and $d_q = d \bmod (q-1)$, computes $m_1 = c^{d_p} \bmod p$ and $m_2 = c^{d_q} \bmod q$, then recombines via Garner's formula.

\textbf{Montgomery Multiplication.}
Montgomery~\cite{montgomery1985} introduced a method for performing modular multiplication without trial division, replacing expensive modulo operations with bitwise shifts. The technique converts operands into ``Montgomery form'' using an auxiliary modulus $R = 2^k$, performs all arithmetic in this domain, and converts back only at the end. This approach is particularly effective in hardware and compiled-language implementations where shift operations are essentially free.

\textbf{Sliding-Window and $k$-ary Exponentiation.}
The standard binary square-and-multiply method processes the exponent one bit at a time, requiring $n$ squarings and approximately $n/2$ multiplications for an $n$-bit exponent. Knuth~\cite{knuth1997} describes the $m$-ary method that processes multiple bits simultaneously. The sliding-window variant~\cite{menezes2001} further reduces the number of multiplications by pre-computing a table of odd powers and scanning the exponent in variable-length windows, reducing multiplications to approximately $n/(w+1) + 2^{w-1} - 1$ for window size $w$.

\textbf{Key Caching and Batch Processing.}
In real-world deployments, RSA key pairs are typically generated once and reused across many sessions. Fiat~\cite{fiat1997} proposed batch RSA, which amortizes the cost of a single modular exponentiation across multiple messages. While full batch RSA requires specific structural properties, the simpler optimization of caching pre-generated keys and reusing them eliminates the dominant key-generation bottleneck in web server scenarios.

\subsection{Cryptanalysis and Security Considerations}

Boneh~\cite{boneh1999} provides a comprehensive survey of attacks on RSA, categorizing them into elementary attacks (small $e$, small $d$), factoring-based attacks, and implementation attacks. Wiener~\cite{wiener1990} showed that RSA with private exponent $d < n^{0.25}$ can be broken efficiently using continued fraction expansion. Coppersmith~\cite{coppersmith1997} developed lattice-based methods for finding small roots of modular polynomial equations, enabling partial key exposure attacks.

\textbf{Implementation Attacks.}
Kocher~\cite{kocher1996} introduced timing attacks, demonstrating that variations in computation time can leak information about secret keys. Kocher, Jaffe, and Jun~\cite{kocher1999} subsequently introduced power analysis attacks. Boneh, DeMillo, and Lipton~\cite{boneh2001} showed that computational faults induced during RSA operations can reveal private keys, motivating fault-attack protection mechanisms in optimized implementations.

\textbf{Protocol-Level Attacks.}
Bleichenbacher~\cite{bleichenbacher1998} demonstrated a practical adaptive chosen-ciphertext attack against RSA encryption using PKCS\#1 v1.5 padding, leading to the adoption of OAEP~\cite{bellare1995}. Manger~\cite{manger2001} extended these results to RSAES-OAEP. Heninger et al.~\cite{heninger2012} conducted a large-scale analysis of public keys on the Internet, finding that a significant fraction of RSA keys shared prime factors due to poor random number generation.

\subsection{Key Size Recommendations}

Lenstra and Verheul~\cite{lenstra2001} developed a systematic framework for selecting cryptographic key sizes based on projected advances in computing power and cryptanalysis. NIST SP 800-57~\cite{nist_sp80057} recommends minimum 2048-bit RSA keys through 2030, and 3072-bit or larger for longer-term protection. Håstad and N\"aslund~\cite{hastad1989} showed that RSA with small exponents can be broken when the same message is encrypted under multiple keys.

\subsection{Hybrid Encryption Architectures}

The hybrid encryption paradigm---using asymmetric cryptography for key exchange and symmetric cryptography for bulk data encryption---is described in the Handbook of Applied Cryptography~\cite{menezes2001} and formalized in PKCS\#1~\cite{pkcs1_2016}. Modern implementations typically use RSA to wrap a random AES-256-GCM key~\cite{nist_fips197}, which then encrypts the actual data. This approach combines RSA's key management capabilities with AES's hardware-accelerated throughput.

\subsection{Post-Quantum Cryptography}

Shor's algorithm~\cite{shor1997} can factor integers and compute discrete logarithms in polynomial time on a quantum computer, rendering RSA and ECC vulnerable. Bernstein~\cite{bernstein2010} analyzed the impact of Grover's algorithm on symmetric cryptography, recommending AES-256 for long-term security. NIST's Post-Quantum Cryptography standardization process~\cite{nist_pqc} has selected CRYSTALS-Kyber~\cite{kyber2023} for key encapsulation and CRYSTALS-Dilithium~\cite{dilithium2022} for digital signatures, both based on the Learning With Errors (LWE) problem over module lattices. Alagic et al.~\cite{alagic2022} provide a comprehensive status report on the third round of NIST's PQC standardization.


\section{Proposed System and Flow of Work}
\label{sec:proposed}

\subsection{Research Methodology}

This study follows a quantitative experimental methodology comprising three phases:

\begin{enumerate}
    \item \textbf{Literature Synthesis:} Systematic review of RSA optimization techniques from 30~key publications spanning 1959--2024, extracting algorithm descriptions, theoretical complexity claims, and reported performance improvements.

    \item \textbf{Controlled Implementation:} Development of two Python-based RSA implementations---a standard (textbook) baseline and an optimized variant---sharing the same mathematical foundation but differing in computational strategy. The implementations serve as controlled testbeds where each optimization can be individually and collectively evaluated.

    \item \textbf{Empirical Benchmarking:} Execution of six benchmark suites measuring correctness, speed, memory consumption, hybrid encryption efficiency, key-size scaling, and concurrent throughput against fresh experimental data.
\end{enumerate}

\subsection{System Architecture}

The experimental system is organized into four modules:

\begin{itemize}
    \item \textbf{standard\_rsa/}: A baseline RSA implementation with no algorithmic optimizations. Implements key generation via Miller-Rabin primality testing~\cite{menezes2001}, textbook encryption/decryption using $c = m^e \bmod n$ and $m = c^d \bmod n$, naive square-and-multiply modular exponentiation, PKCS\#1 v1.5-style padding~\cite{pkcs1_2016}, and hybrid RSA+AES-GCM. Generates fresh keys per call with no caching.

    \item \textbf{optimized\_rsa/}: An enhanced RSA implementation incorporating eight optimizations from the surveyed literature: (1)~CRT-based decryption~\cite{quisquater1982}; (2)~sliding-window modular exponentiation with 4-bit window~\cite{knuth1997}; (3)~Montgomery multiplication~\cite{montgomery1985}; (4)~key caching via a \texttt{KeyCache} class; (5)~batch encrypt/decrypt functions; (6)~pre-computed CRT parameters ($d_p$, $d_q$, $q_{\mathrm{inv}}$); (7)~hybrid RSA+AES-256-GCM envelope encryption; (8)~fault-attack verification~\cite{boneh2001}.

    \item \textbf{tests/}: Six benchmark test files covering correctness verification, speed measurement, memory profiling, hybrid encryption comparison, key-size scaling, and concurrency simulation.

    \item \textbf{perf\_eval.py}: The primary benchmark runner using the \texttt{cryptography} library~\cite{cryptography_io} for AES-256-GCM operations and \texttt{tracemalloc}~\cite{python_tracemalloc} for memory profiling.
\end{itemize}

\subsection{Flow of Work}

The experimental workflow proceeds as follows:
\begin{enumerate}
    \item Run correctness verification (40~assertions) to validate both implementations.
    \item Execute speed benchmarks (50~samples per measurement) for encrypt/decrypt operations.
    \item Measure memory allocation using \texttt{tracemalloc} across all operations.
    \item Compare hybrid RSA+AES-GCM versus pure RSA chunked encryption across six data sizes (1~KB to 1~MB).
    \item Evaluate key-size scaling across 1024, 2048, 3072, and 4096-bit keys.
    \item Simulate concurrent web sessions using \texttt{ThreadPoolExecutor} with 1--50 concurrent threads.
    \item Synthesize results with the literature survey to produce a comparative analysis.
\end{enumerate}


\section{Algorithm and Comparative Study}
\label{sec:comparative}

This section presents a detailed comparison of the RSA optimization algorithms identified in the literature survey, organized by their theoretical properties, implementation complexity, and expected practical impact.

\subsection{Comparative Overview of Optimization Techniques}

Table~\ref{tab:algo_compare} provides a side-by-side comparison of the major RSA optimization techniques surveyed.

\begin{table*}[!t]
\caption{Comparative Analysis of RSA Optimization Techniques from Literature}
\label{tab:algo_compare}
\centering
\small
\begin{tabular}{|p{2.2cm}|p{1.5cm}|p{2.8cm}|p{2.2cm}|p{1.5cm}|p{2.5cm}|}
\hline
\textbf{Technique} & \textbf{Origin} & \textbf{Mechanism} & \textbf{Theoretical Speedup} & \textbf{Applies To} & \textbf{Limitation} \\
\hline
CRT Decryption & Quisquater \& Couvreur~\cite{quisquater1982} & Split $c^d \bmod n$ into two half-size exponentiations via Garner's formula & $3$--$4\times$ & Decryption only & Vulnerable to Bellcore fault attack~\cite{boneh2001}; requires fault verification \\
\hline
Montgomery Multiplication & Montgomery~\cite{montgomery1985} & Replace modular reduction (division) with shift-based REDC algorithm & $2$--$3\times$ per multiplication & All modular arithmetic & Conversion overhead; less effective in interpreted languages \\
\hline
Sliding-Window Exponentiation & Knuth~\cite{knuth1997}, Menezes et al.~\cite{menezes2001} & Pre-compute odd power table; scan exponent in $w$-bit windows & Reduces multiplications by $30$--$40\%$ & All exponentiations & Table storage overhead; marginal gain for small exponents \\
\hline
Key Caching & Practical deployment~\cite{fiat1997} & Generate keys once, reuse across sessions & $40$--$2000\times$ (vs.\ cold start) & Full encrypt/decrypt cycle & Key rotation complexity; forward secrecy trade-off \\
\hline
Batch Processing & Fiat~\cite{fiat1997} & Amortize single exponentiation across $k$ messages & $O(\sqrt{k})$ amortized cost per message & Multiple messages & Requires structural pre-conditions; limited batch size \\
\hline
Hybrid RSA+AES & Envelope encryption~\cite{menezes2001}, PKCS\#1~\cite{pkcs1_2016} & RSA wraps AES-256 key; AES encrypts bulk data & $8$--$760\times$ for bulk data & Data $>$ 256 bytes & Two-algorithm complexity; key management overhead \\
\hline
\end{tabular}
\end{table*}

\subsection{Complexity Analysis}

The computational complexity of RSA's core operation---modular exponentiation $a^b \bmod n$---depends on the method used:

\textbf{Binary square-and-multiply (standard):}
Requires $k$ squarings and up to $k$ multiplications for a $k$-bit exponent, each multiplication being $O(\ell^2)$ for $\ell$-bit operands, yielding $O(k \cdot \ell^2)$ total complexity.

\textbf{CRT-based decryption:}
Replaces one exponentiation with $k$-bit operands by two exponentiations with $k/2$-bit operands:
\begin{align}
    m_1 &= c^{d_p} \bmod p \label{eq:crt1} \\
    m_2 &= c^{d_q} \bmod q \label{eq:crt2} \\
    h &= q_{\mathrm{inv}} \cdot (m_1 - m_2) \bmod p \label{eq:crt3} \\
    m &= m_2 + h \cdot q \label{eq:crt4}
\end{align}
Total complexity: $2 \cdot O(k \cdot (\ell/2)^2) = O(k \cdot \ell^2 / 2)$, a theoretical $4\times$ improvement (reduced by recombination overhead to approximately $3\times$).

\textbf{Sliding-window method:}
Pre-computes a table of $2^{w-1}$ odd powers and processes the exponent in windows of up to $w$ bits. The average number of multiplications drops from $k/2$ to approximately $k/(w+1) + 2^{w-1}$, which for $w=4$ and $k=2048$ yields approximately 425 multiplications versus 1024 for binary method.

\textbf{Montgomery multiplication:}
Each modular multiplication $a \cdot b \bmod n$ replaces the expensive division by $n$ with $O(\ell)$ additions and shifts via the REDC algorithm. In hardware, this provides $2$--$3\times$ speedup per multiplication. However, the conversion to/from Montgomery form adds overhead of $O(\ell^2)$ at the beginning and end of each exponentiation.

\subsection{Comparative Study of Surveyed Papers}

Table~\ref{tab:papers_compare} summarizes the key contributions from the surveyed literature and their reported findings.

\begin{table*}[!t]
\caption{Summary of Key Publications in RSA Optimization and Security}
\label{tab:papers_compare}
\centering
\small
\begin{tabular}{|p{3.5cm}|p{1.5cm}|p{3cm}|p{4.5cm}|}
\hline
\textbf{Contribution} & \textbf{Year} & \textbf{Authors} & \textbf{Key Finding} \\
\hline
Public-key cryptography concept & 1976 & Diffie \& Hellman~\cite{diffie1976} & Introduced the concept of asymmetric encryption and key exchange \\
\hline
RSA cryptosystem & 1978 & Rivest, Shamir, Adleman~\cite{rivest1978} & First practical public-key encryption and signature scheme \\
\hline
CRT-accelerated RSA & 1982 & Quisquater \& Couvreur~\cite{quisquater1982} & CRT yields $\sim$$4\times$ speedup for RSA decryption \\
\hline
Montgomery multiplication & 1985 & Montgomery~\cite{montgomery1985} & Modular multiplication without trial division \\
\hline
Elliptic curve cryptography & 1985/1987 & Miller~\cite{miller1985}, Koblitz~\cite{koblitz1987} & Equivalent security to RSA with much smaller keys \\
\hline
Timing attacks on RSA & 1996 & Kocher~\cite{kocher1996} & Execution time variations leak secret key information \\
\hline
Quantum factoring algorithm & 1997 & Shor~\cite{shor1997} & Polynomial-time integer factorization on quantum computers \\
\hline
PKCS\#1 v1.5 attack & 1998 & Bleichenbacher~\cite{bleichenbacher1998} & Adaptive chosen-ciphertext attack on RSA with PKCS\#1 padding \\
\hline
Survey of RSA attacks & 1999 & Boneh~\cite{boneh1999} & Comprehensive taxonomy of attacks on the RSA cryptosystem \\
\hline
Fault attacks on RSA & 2001 & Boneh, DeMillo, Lipton~\cite{boneh2001} & Hardware faults during CRT computation can reveal private keys \\
\hline
Key size selection framework & 2001 & Lenstra \& Verheul~\cite{lenstra2001} & Systematic method for choosing cryptographic key lengths \\
\hline
Weak RSA keys on the Internet & 2012 & Heninger et al.~\cite{heninger2012} & $0.2\%$ of RSA public keys share a prime factor due to poor RNG \\
\hline
CRYSTALS-Kyber (PQC) & 2023 & Avanzi et al.~\cite{kyber2023} & Lattice-based KEM resistant to classical and quantum attacks \\
\hline
CRYSTALS-Dilithium (PQC) & 2022 & Bai et al.~\cite{dilithium2022} & Lattice-based digital signature scheme for post-quantum era \\
\hline
NIST PQC standardization & 2022 & Alagic et al.~\cite{alagic2022} & Third-round status report on PQC standardization process \\
\hline
\end{tabular}
\end{table*}


\section{Experimental Results}
\label{sec:results}

\subsection{Test Environment}

All experiments were conducted on the following platform:
\begin{itemize}
    \item \textbf{Processor:} Apple Silicon (ARM64), 10 performance cores
    \item \textbf{Memory:} 16~GB unified memory
    \item \textbf{Operating System:} macOS 26.4.1 (Darwin 25.4.0)
    \item \textbf{Python:} 3.14.4 with \texttt{cryptography} 46.0.7~\cite{cryptography_io}
    \item \textbf{Memory Profiler:} \texttt{tracemalloc}~\cite{python_tracemalloc}
    \item \textbf{Concurrency:} \texttt{threading} + \texttt{concurrent.futures}~\cite{python_threading}
\end{itemize}

All RSA operations were implemented from scratch; AES-256-GCM operations use the \texttt{cryptography} library's C-level implementation for realistic hybrid encryption benchmarks.

\newpage

\subsection{Implementation Details}

The optimized RSA implementation comprises eight integrated modules. The key source code excerpts are presented below.

\textbf{Primality Testing (Miller-Rabin).}
Candidates are trial-divided against the first 27~primes before $k=20$ rounds of Miller-Rabin testing. The most-significant and least-significant bits are forced to~1 to ensure correct bit length and oddness.

\begin{lstlisting}
_SMALL_PRIMES = [2,3,5,7,11,13,17,19,23,29,31,37,
                 41,43,47,53,59,61,67,71,73,79,
                 83,89,97]

def is_probable_prime(n: int, k: int = 20) -> bool:
    if n < 2: return False
    for p in _SMALL_PRIMES:
        if n % p == 0: return n == p
    r, d = 0, n - 1
    while d % 2 == 0: d //= 2; r += 1
    for _ in range(k):
        a = random.randrange(2, n - 1)
        x = pow(a, d, n)
        if x == 1 or x == n - 1: continue
        for _ in range(r - 1):
            x = pow(x, 2, n)
            if x == n - 1: break
        else: return False
    return True

def generate_prime(bits: int) -> int:
    while True:
        candidate = random.getrandbits(bits)
        candidate |= (1 << (bits - 1)) | 1
        if is_probable_prime(candidate):
            return candidate
\end{lstlisting}
\vspace{-2mm}
\captionof{lstlisting}{Miller-Rabin primality test and prime generation}
\label{lst:prime}

\textbf{Key Generation with CRT Parameters.}
The private key is extended with pre-computed CRT parameters $d_p$, $d_q$, and $q_{\mathrm{inv}}$ for accelerated decryption.

\begin{lstlisting}
def generate_keypair(key_size: int = 2048):
    half = key_size // 2
    p = generate_prime(half)
    q = generate_prime(half)
    while q == p:
        q = generate_prime(half)
    n = p * q
    phi = (p - 1) * (q - 1)
    e = 65537
    d = pow(e, -1, phi)
    # CRT parameters
    dp = d % (p - 1)
    dq = d % (q - 1)
    qinv = pow(q, -1, p)
    pub = {"e": e, "n": n}
    priv = {"d": d, "n": n, "p": p, "q": q,
            "dp": dp, "dq": dq, "qinv": qinv}
    return priv, pub
\end{lstlisting}
\vspace{-2mm}
\captionof{lstlisting}{RSA keypair generation with CRT parameters}
\label{lst:keygen}

\textbf{Montgomery Multiplication.}
The \texttt{MontgomeryContext} class implements the REDC algorithm, replacing expensive modular division with bitwise shifts.

\begin{lstlisting}
class MontgomeryContext:
    def __init__(self, n: int):
        self.n = n
        self.r = 1 << (n.bit_length())
        self.r_inv = pow(self.r, -1, n)
        self.n_prime = (-pow(n,-1,self.r)) % self.r
        self.r2 = (self.r * self.r) % n

    def to_mont(self, x: int) -> int:
        return (x * self.r) % self.n

    def from_mont(self, x: int) -> int:
        return (x * self.r_inv) % self.n

    def reduce(self, t: int) -> int:
        n, r, n_prime = self.n, self.r, self.n_prime
        m = (t * n_prime) % r
        t = (t + m * n) // r
        if t >= n: t -= n
        return t

    def mul(self, a: int, b: int) -> int:
        return self.reduce(a * b)
\end{lstlisting}
\vspace{-2mm}
\captionof{lstlisting}{Montgomery multiplication (REDC)}
\label{lst:montgomery}

\textbf{Sliding-Window Exponentiation.}
A variable window size ($w=4$ for large exponents, $w=2$ for small) with pre-computed odd power table reduces multiplications from $k/2$ to approximately $k/(w+1) + 2^{w-1}$.

\begin{lstlisting}
def mod_exp_sliding_window(base, exp, mod):
    if exp == 0: return 1 % mod
    if mod == 1: return 0
    if mod % 2 == 0:
        return sliding_window_no_mont(base, exp, mod)

    mont = MontgomeryContext(mod)
    w = 4 if exp.bit_length() > 64 else 2
    g = mont.to_mont(base)
    g2 = mont.mul(g, g)
    table = {1: g}
    for i in range(3, 1 << w, 2):
        table[i] = mont.mul(table[i - 2], g2)

    result = mont.to_mont(1)
    i = exp.bit_length() - 1
    while i >= 0:
        if not (exp >> i) & 1:
            result = mont.mul(result, result)
            i -= 1
        else:
            best_j, best_val = i, 1
            j = i
            while j >= 0 and (i - j + 1) <= w:
                val = (exp>>j) & ((1<<(i-j+1))-1)
                if val & 1: best_j, best_val = j, val
                j -= 1
            for _ in range(i - best_j + 1):
                result = mont.mul(result, result)
            result = mont.mul(result, table[best_val])
            i = best_j - 1
    return mont.from_mont(result)
\end{lstlisting}
\vspace{-2mm}
\captionof{lstlisting}{Sliding-window modular exponentiation with Montgomery form}
\label{lst:sliding}

\textbf{CRT-Based Decryption with Fault Verification.}
Decomposes decryption into two half-size exponentiations via CRT and verifies the result against the Bellcore fault attack~\cite{boneh2001}.

\begin{lstlisting}
def decrypt(priv: dict, ciphertext: int) -> int:
    p, q = priv["p"], priv["q"]
    # CRT half-size exponentiations
    m1 = mod_exp_sliding_window(ciphertext, priv["dp"], p)
    m2 = mod_exp_sliding_window(ciphertext, priv["dq"], q)
    # Garner's recombination
    h = (priv["qinv"] * (m1 - m2)) % p
    m = m2 + h * q
    # Fault-attack check
    if pow(m, 65537, priv["n"]) != ciphertext % priv["n"]:
        raise RuntimeError("Fault detected in CRT decryption")
    return m
\end{lstlisting}
\vspace{-2mm}
\captionof{lstlisting}{CRT-based decryption with fault-attack protection}
\label{lst:decrypt}

\textbf{PKCS\#1 v1.5 Padding.}
Implements PKCS\#1 v1.5-style padding, limiting the payload to 245~bytes for 2048-bit keys.

\begin{lstlisting}
def pkcs1_pad(data: bytes, n_bytes: int) -> int:
    pad_len = n_bytes - len(data) - 3
    if pad_len < 8:
        raise ValueError("Data too long for key size")
    padding = bytes(random.randint(1, 255)
                    for _ in range(pad_len))
    padded = b"\x00\x02" + padding + b"\x00" + data
    return int.from_bytes(padded, "big")

def pkcs1_unpad(padded_int: int) -> bytes:
    data = padded_int.to_bytes(
        (padded_int.bit_length()+7)//8, "big")
    idx = data.index(b"\x00", 2)
    return data[idx + 1:]
\end{lstlisting}
\vspace{-2mm}
\captionof{lstlisting}{PKCS\#1 v1.5 padding and unpadding}
\label{lst:padding}

\textbf{Key Cache.}
The \texttt{KeyCache} class lazily generates a keypair on first access and reuses it for all subsequent operations, eliminating the dominant key-generation bottleneck.

\begin{lstlisting}
class KeyCache:
    def __init__(self, key_size: int = 2048):
        self.key_size = key_size
        self._priv = None
        self._pub = None

    def get(self):
        if self._priv is None:
            self._priv, self._pub = \
                generate_keypair(self.key_size)
        return self._priv, self._pub

    def regenerate(self):
        self._priv, self._pub = \
            generate_keypair(self.key_size)
        return self._priv, self._pub
\end{lstlisting}
\vspace{-2mm}
\captionof{lstlisting}{Key cache for reuse across sessions}
\label{lst:keycache}

\textbf{Hybrid RSA+AES-256-GCM.}
RSA wraps a random AES-256 key; AES-GCM encrypts bulk data. The envelope stores four fields: wrapped key, IV, ciphertext, and authentication tag.

\begin{lstlisting}
from cryptography.hazmat.primitives.ciphers \
    import Cipher, algorithms, modes

def hybrid_encrypt(pub: dict, data: bytes) -> dict:
    aes_key = os.urandom(32)
    iv = os.urandom(12)
    enc = Cipher(algorithms.AES(aes_key),
                 modes.GCM(iv)).encryptor()
    ct = enc.update(data) + enc.finalize()
    wrapped_key = encrypt_bytes(pub, aes_key)
    return {"wrapped_key": wrapped_key, "iv": iv,
            "ct": ct, "tag": enc.tag}

def hybrid_decrypt(priv: dict, envelope: dict) -> bytes:
    aes_key = decrypt_bytes(priv, envelope["wrapped_key"])
    dec = Cipher(algorithms.AES(aes_key),
        modes.GCM(envelope["iv"],
                  envelope["tag"])).decryptor()
    return dec.update(envelope["ct"]) + dec.finalize()
\end{lstlisting}
\vspace{-2mm}
\captionof{lstlisting}{Hybrid RSA+AES-256-GCM envelope encryption}
\label{lst:hybrid}

\subsection{Standard vs.\ Optimized RSA (2048-bit)}

Table~\ref{tab:std_vs_opt} presents the direct comparison between standard and optimized RSA implementations using 2048-bit keys.

\begin{table}[H]
\caption{Standard vs.\ Optimized RSA Performance (2048-bit, from-scratch)}
\label{tab:std_vs_opt}
\centering
\begin{tabular}{|c|c|c|c|}
\hline
\textbf{Operation} & \textbf{Standard (ms)} & \textbf{Optimized (ms)} & \textbf{Speedup} \\
\hline
Encrypt & 2910.4 & 2.6 & $1112\times$ \\
\hline
Decrypt & 3142.1 & 25.6 & $123\times$ \\
\hline
Enc Memory (KB) & 5.4 & 5.4 & $1.0\times$ \\
\hline
Dec Memory (KB) & 4.3 & 5.1 & $0.8\times$ \\
\hline
\end{tabular}
\end{table}

The dramatic speedup is primarily attributable to key caching in the optimized implementation. The standard implementation generates fresh keys per operation (cold start), while the optimized variant reuses cached keys. The encryption speedup of $1112\times$ reflects the elimination of the 2048-bit key generation bottleneck (which dominates the standard implementation's 2910~ms encrypt time, as pure Python Miller-Rabin primality testing is inherently slow). Decryption benefits from both key caching and CRT-based computation, yielding a $123\times$ improvement. Memory usage is comparable between implementations, with the optimized variant using slightly more for decryption due to CRT parameter storage.

\subsection{Key-Size Scaling (1024--4096 bits)}

Table~\ref{tab:keysize} shows how the optimized RSA implementation scales across four key sizes.

\begin{table}[H]
\caption{Key-Size Scaling Performance (Optimized RSA, from-scratch, 30 Samples)}
\label{tab:keysize}
\centering
\begin{tabular}{|c|c|c|c|c|}
\hline
\textbf{Key Size} & \textbf{Enc (ms)} & \textbf{Dec (ms)} & \textbf{Enc Mem (KB)} & \textbf{Dec Mem (KB)} \\
\hline
1024-bit & 1.082 & 7.495 & 2.90 & 3.07 \\
\hline
2048-bit & 2.605 & 25.735 & 5.27 & 5.21 \\
\hline
3072-bit & 4.740 & 69.446 & 9.47 & 7.34 \\
\hline
4096-bit & 7.134 & 153.323 & 12.36 & 9.20 \\
\hline
\end{tabular}
\end{table}

Encryption time grows from 1.082~ms (1024-bit) to 7.134~ms (4096-bit), a $6.6\times$ increase for a $4\times$ key-size increase. This sub-quadratic growth occurs because the public exponent $e = 65537$ remains fixed at 17~bits regardless of key size. Decryption time grows steeply from 7.495~ms to 153.323~ms, a $20.5\times$ increase, consistent with the $O(n^2 \log n)$ complexity of modular exponentiation in pure Python.

The growth factors per key-size increment are:
\begin{itemize}
    \item 1024 $\rightarrow$ 2048: encrypt $2.4\times$, decrypt $3.4\times$
    \item 2048 $\rightarrow$ 3072: encrypt $1.8\times$, decrypt $2.7\times$
    \item 3072 $\rightarrow$ 4096: encrypt $1.5\times$, decrypt $2.2\times$
\end{itemize}

\subsection{Hybrid RSA+AES-GCM vs.\ Pure RSA}

Table~\ref{tab:hybrid} confirms the overwhelming superiority of the hybrid approach for bulk data encryption.

\begin{table}[H]
\caption{Hybrid RSA+AES-GCM vs.\ Pure RSA Encryption Time (from-scratch)}
\label{tab:hybrid}
\centering
\begin{tabular}{|c|c|c|c|}
\hline
\textbf{Data Size} & \textbf{Pure RSA (ms)} & \textbf{Hybrid (ms)} & \textbf{Speedup} \\
\hline
1~KB & 13.95 & 3.06 & $5\times$ \\
\hline
10~KB & 117.99 & 2.84 & $42\times$ \\
\hline
50~KB & 587.67 & 2.86 & $206\times$ \\
\hline
100~KB & 1174.63 & 2.88 & $408\times$ \\
\hline
500~KB & 5867.58 & 2.91 & $2015\times$ \\
\hline
1000~KB & 11768.01 & 3.02 & $3903\times$ \\
\hline
\end{tabular}
\end{table}

Pure RSA encrypts data in 245-byte chunks, requiring multiple RSA operations proportional to data size. The hybrid approach performs exactly one RSA operation (to wrap the 32-byte AES key) regardless of data size, with AES-GCM handling bulk encryption at near-constant time. At 1~MB, the hybrid approach is $3903\times$ faster. The hybrid encryption time remains nearly constant from 2.84~ms at 10~KB to 3.02~ms at 1~MB, confirming that AES-GCM throughput dominates.

\subsection{Concurrent Throughput}

Table~\ref{tab:concurrency} presents the dynamic load simulation results.

\begin{table}[H]
\caption{Concurrent Session Throughput (Optimized RSA, 2048-bit, from-scratch)}
\label{tab:concurrency}
\centering
\begin{tabular}{|c|c|c|c|}
\hline
\textbf{Threads} & \textbf{Throughput (sess/s)} & \textbf{Avg Latency (ms)} & \textbf{Peak Mem (MB)} \\
\hline
1 & 34.9 & 28.6 & 0.11 \\
\hline
5 & 34.5 & 136.7 & 0.12 \\
\hline
10 & 34.1 & 250.3 & 0.15 \\
\hline
25 & 34.2 & 254.7 & 0.21 \\
\hline
50 & 34.2 & 355.6 & 0.22 \\
\hline
\end{tabular}
\end{table}

The optimized implementation maintains a stable throughput of approximately 34--35~sessions/second across all concurrency levels. This near-constant throughput reflects Python's Global Interpreter Lock (GIL), which serializes CPU-bound work across threads. Average latency increases linearly with concurrency (from 28.6~ms at 1~thread to 355.6~ms at 50~threads) as threads queue for the GIL. Memory usage grows modestly from 0.11~MB to 0.22~MB. While the per-thread throughput is lower than compiled implementations would achieve, the key caching strategy ensures that each session completes in a consistent time without key-generation overhead.

\subsection{Summary of Comparative Findings}

Table~\ref{tab:summary} synthesizes the experimental results with the theoretical predictions from the literature.

\begin{table}[H]
\caption{Summary: Theoretical vs.\ Experimental Speedup (from-scratch, 2048-bit)}
\label{tab:summary}
\centering
\begin{tabular}{|c|c|c|}
\hline
\textbf{Optimization} & \textbf{Theory} & \textbf{Observed} \\
\hline
CRT Decryption & $3$--$4\times$ & $1.27\times$ (2048-bit) \\
\hline
Montgomery Mult. & $2$--$3\times$ & $0.33\times$ (Python) \\
\hline
Sliding-Window & $30$--$40\%$ fewer mult. & $1.28\times$ (Python) \\
\hline
Key Caching & $N/A$ (deployment) & $1112$--$2117\times$ \\
\hline
Hybrid RSA+AES & $8$--$400\times$ & $5$--$3903\times$ \\
\hline
\end{tabular}
\end{table}

The discrepancy between theoretical and observed performance for Montgomery multiplication and sliding-window exponentiation is attributable to Python's interpreter overhead: per-operation costs (object creation, function dispatch, garbage collection) dominate the actual arithmetic for moderate key sizes. In particular, Montgomery multiplication is $3\times$ \textit{slower} than standard arithmetic in Python ($69.3$~ms vs.\ $22.9$~ms for 2048-bit modular exponentiation) because the conversion to/from Montgomery form and the REDC algorithm's additional integer operations incur Python interpreter overhead that outweighs the savings from eliminating trial division. Key caching and hybrid encryption, which reduce the \textit{number} of expensive operations rather than per-operation cost, deliver the largest practical improvements.


\section{Conclusion}
\label{sec:conclusion}

This paper presented a comparative study of RSA optimization techniques for web application security, synthesizing findings from 30~key publications and validating them through controlled experiments. The study surveyed six major optimization approaches---CRT-based decryption, Montgomery multiplication, sliding-window exponentiation, key caching, batch processing, and hybrid RSA+AES encryption---and evaluated their practical impact through fresh benchmarks across six dimensions.

The principal findings are:

\begin{enumerate}
    \item \textbf{Key caching is the highest-impact optimization}, providing $1112$--$2117\times$ speedup by eliminating repeated key generation. This validates the importance of key management infrastructure in web server deployments.

    \item \textbf{Hybrid RSA+AES-GCM is essential for bulk data}, outperforming pure RSA by $5\times$ to $3903\times$ across 1~KB to 1~MB data sizes. RSA should never be used alone for data exceeding a few hundred bytes.

    \item \textbf{CRT provides measurable decryption improvement} ($1.27\times$ at 2048-bit), though the theoretical $4\times$ prediction assumes compiled-language execution. In interpreted environments, the gain is modest.

    \item \textbf{Montgomery multiplication and sliding-window exponentiation underperform in Python} due to interpreter overhead. Montgomery multiplication is actually $3\times$ slower in pure Python ($69.3$~ms vs.\ $22.9$~ms). These techniques are better suited to C, Rust, or hardware implementations.

    \item \textbf{RSA scales approximately quadratically with key size}, with decryption growing from 7.5~ms (1024-bit) to 153.3~ms (4096-bit), a $20.5\times$ increase. This makes the transition to 3072-bit or 4096-bit keys (recommended by NIST~\cite{nist_sp80057}) a significant performance concern that must be mitigated through the optimizations surveyed.

    \item \textbf{Concurrent throughput is bounded by Python's GIL} at approximately 34~sessions/second regardless of thread count. While key caching eliminates the key-generation bottleneck, true parallelism requires compiled implementations or multiprocessing.
\end{enumerate}

These findings support three recommendations for web application architects: (1)~adopt hybrid encryption as the default for all data payloads, (2)~implement key caching with periodic rotation to eliminate the key-generation bottleneck, and (3)~begin planning for post-quantum migration to lattice-based schemes~\cite{kyber2023}, \cite{dilithium2022} as RSA's security margin narrows. Future work should replicate these benchmarks in compiled languages, evaluate energy consumption for edge deployments, and compare RSA directly against CRYSTALS-Kyber in hybrid configurations.


\section*{Acknowledgment}

The authors would like to thank the faculty of the Department of Information Technology, C.~K.~Pithawala College of Engineering, for their guidance and support throughout this research.


\ifCLASSOPTIONcaptionsoff
\newpage
\fi

\bibliographystyle{IEEEtran}
\begin{thebibliography}{30}

\bibitem{rivest1978}
R.~L. Rivest, A.~Shamir, and L.~Adleman, ``A method for obtaining digital signatures and public-key cryptosystems,'' \textit{Commun. ACM}, vol.~21, no.~2, pp.~120--126, 1978. [Online]. Available: \url{https://doi.org/10.1145/359340.359342}

\bibitem{boneh1999}
D.~Boneh, ``Twenty years of attacks on the RSA cryptosystem,'' \textit{Notices Amer. Math. Soc.}, vol.~46, no.~2, pp.~203--213, 1999. [Online]. Available: \url{https://www.ams.org/notices/199902/boneh.pdf}

\bibitem{coppersmith1997}
D.~Coppersmith, ``Small solutions to polynomial equations, and low exponent RSA vulnerabilities,'' \textit{J. Cryptol.}, vol.~10, no.~4, pp.~233--260, 1997. [Online]. Available: \url{https://doi.org/10.1007/s001459900030}

\bibitem{shor1997}
P.~W. Shor, ``Polynomial-time algorithms for prime factorization and discrete logarithms on a quantum computer,'' \textit{SIAM J. Comput.}, vol.~26, no.~5, pp.~1484--1509, 1997. [Online]. Available: \url{https://doi.org/10.1137/S0097539795293172}

\bibitem{wiener1990}
M.~Wiener, ``Cryptanalysis of short RSA secret exponents,'' \textit{IEEE Trans. Inf. Theory}, vol.~36, no.~3, pp.~553--558, May 1990. [Online]. Available: \url{https://doi.org/10.1109/18.54902}

\bibitem{boneh2001}
D.~Boneh, R.~A. DeMillo, and R.~J. Lipton, ``On the importance of eliminating errors in cryptographic computations,'' \textit{J. Cryptol.}, vol.~14, no.~2, pp.~101--119, 2001. [Online]. Available: \url{https://doi.org/10.1007/s001450010016}

\bibitem{garner1959}
H.~L. Garner, ``The residue number system,'' \textit{IRE Trans. Electron. Comput.}, vol.~EC-8, no.~2, pp.~140--147, Jun. 1959. [Online]. Available: \url{https://doi.org/10.1109/TEC.1959.5219515}

\bibitem{montgomery1985}
P.~L. Montgomery, ``Modular multiplication without trial division,'' \textit{Math. Comput.}, vol.~44, no.~170, pp.~519--521, Apr. 1985. [Online]. Available: \url{https://www.ams.org/journals/mcom/1985-44-170/S0025-5718-1985-0777282-X/}

\bibitem{quisquater1982}
J.-J. Quisquater and C.~Couvreur, ``Fast decipherment algorithm for RSA public-key cryptosystem,'' \textit{Electron. Lett.}, vol.~18, no.~21, pp.~905--907, Oct. 1982. [Online]. Available: \url{https://doi.org/10.1049/el:19820617}

\bibitem{diffie1976}
W.~Diffie and M.~E. Hellman, ``New directions in cryptography,'' \textit{IEEE Trans. Inf. Theory}, vol.~IT-22, no.~6, pp.~644--654, Nov. 1976. [Online]. Available: \url{https://doi.org/10.1109/TIT.1976.1055638}

\bibitem{rabin1979}
M.~O. Rabin, ``Digitalized signatures and public-key functions as intractable as factorization,'' MIT Lab. Comput. Sci., Cambridge, MA, USA, Tech. Rep. MIT/LCS/TR-212, 1979. [Online]. Available: \url{https://hdl.handle.net/1721.1/146846}

\bibitem{elgamal1985}
T.~ElGamal, ``A public key cryptosystem and a signature scheme based on discrete logarithms,'' \textit{IEEE Trans. Inf. Theory}, vol.~IT-31, no.~4, pp.~469--472, Jul. 1985. [Online]. Available: \url{https://doi.org/10.1109/TIT.1985.1057074}

\bibitem{koblitz1987}
N.~Koblitz, ``Elliptic curve cryptosystems,'' \textit{Math. Comput.}, vol.~48, no.~177, pp.~203--209, 1987. [Online]. Available: \url{https://www.ams.org/journals/mcom/1987-48-177/S0025-5718-1987-0866109-5/}

\bibitem{miller1985}
V.~S. Miller, ``Use of elliptic curves in cryptography,'' in \textit{Advances in Cryptology---CRYPTO '85}, LNCS, vol.~218. Springer, 1985, pp.~417--426. [Online]. Available: \url{https://doi.org/10.1007/3-540-39799-X_31}

\bibitem{kocher1996}
P.~C. Kocher, ``Timing attacks on implementations of Diffie-Hellman, RSA, DSS, and other systems,'' in \textit{Advances in Cryptology---CRYPTO '96}, LNCS, vol.~1109. Springer, 1996, pp.~104--113. [Online]. Available: \url{https://doi.org/10.1007/3-540-68697-5_9}

\bibitem{kocher1999}
P.~Kocher, J.~Jaffe, and B.~Jun, ``Differential power analysis,'' in \textit{Advances in Cryptology---CRYPTO '99}, LNCS, vol.~1666. Springer, 1999, pp.~388--397. [Online]. Available: \url{https://doi.org/10.1007/3-540-48405-1_25}

\bibitem{bleichenbacher1998}
D.~Bleichenbacher, ``Chosen ciphertext attacks against protocols based on the RSA encryption standard PKCS \#1,'' in \textit{Advances in Cryptology---CRYPTO '98}, LNCS, vol.~1462. Springer, 1998, pp.~1--12. [Online]. Available: \url{https://doi.org/10.1007/BFb0055716}

\bibitem{manger2001}
J.~Manger, ``A chosen ciphertext attack on RSA optimal asymmetric encryption padding (OAEP) as standardized in PKCS \#1 v2.0,'' in \textit{Advances in Cryptology---CRYPTO 2001}, LNCS, vol.~2139. Springer, 2001, pp.~230--238. [Online]. Available: \url{https://doi.org/10.1007/3-540-44647-8_14}

\bibitem{bellare1995}
M.~Bellare and P.~Rogaway, ``Optimal asymmetric encryption padding,'' in \textit{Advances in Cryptology---EUROCRYPT '94}, LNCS, vol.~950. Springer, 1995, pp.~92--111. [Online]. Available: \url{https://doi.org/10.1007/BFb0053428}

\bibitem{heninger2012}
N.~Heninger, Z.~Durumeric, E.~Wustrow, and J.~A. Halderman, ``Mining your Ps and Qs: Detection of widespread weak keys in network devices,'' in \textit{Proc. 21st USENIX Secur. Symp.}, 2012, pp.~205--220. [Online]. Available: \url{https://www.usenix.org/conference/usenixsecurity12/technical-sessions/presentation/heninger}

\bibitem{lenstra2001}
A.~K. Lenstra and E.~R. Verheul, ``Selecting cryptographic key sizes,'' \textit{J. Cryptol.}, vol.~14, no.~4, pp.~255--293, 2001. [Online]. Available: \url{https://doi.org/10.1007/s00145-001-0009-4}

\bibitem{hastad1989}
J.~H{\aa}stad, ``Some optimal inapproximability results,'' in \textit{Proc. 29th Annu. ACM Symp. Theory Comput.}, 1997, pp.~1--10. (Original RSA-related work: J.~H{\aa}stad, ``Solving simultaneous modular equations of low degree,'' \textit{SIAM J. Comput.}, vol.~17, no.~2, pp.~336--341, 1988.) [Online]. Available: \url{https://doi.org/10.1137/0217019}

\bibitem{may2003}
A.~May, ``New partial key exposure attacks on RSA,'' in \textit{Advances in Cryptology---CRYPTO 2003}, LNCS, vol.~2729. Springer, 2003, pp.~27--43. [Online]. Available: \url{https://doi.org/10.1007/978-3-540-45146-4_2}

\bibitem{knuth1997}
D.~E. Knuth, \textit{The Art of Computer Programming, Vol.~2: Seminumerical Algorithms}, 3rd~ed.\ Reading, MA, USA: Addison-Wesley, 1997.

\bibitem{fiat1997}
A.~Fiat, ``Batch RSA,'' in \textit{Advances in Cryptology---CRYPTO '89}, LNCS, vol.~435. Springer, 1990, pp.~175--185. [Online]. Available: \url{https://doi.org/10.1007/0-387-34805-0_17}

\bibitem{pkcs1_2016}
K.~Moriarty, B.~Kaliski, J.~Jonsson, and A.~Ruschka, ``PKCS \#1: RSA cryptography specifications version 2.2,'' RFC 8017, IETF, Nov. 2016. [Online]. Available: \url{https://doi.org/10.17487/RFC8017}

\bibitem{nist_fips197}
National Institute of Standards and Technology, ``Advanced Encryption Standard (AES),'' FIPS PUB 197, NIST, Gaithersburg, MD, USA, Nov. 2001. [Online]. Available: \url{https://csrc.nist.gov/publications/detail/fips/197/final}

\bibitem{nist_sp80057}
E.~Barker, ``Recommendation for key management: Part 1---General,'' NIST SP 800-57 Part 1 Rev.~5, NIST, Gaithersburg, MD, USA, May 2020. [Online]. Available: \url{https://doi.org/10.6028/NIST.SP.800-57pt1r5}

\bibitem{nist_pqc}
National Institute of Standards and Technology, ``Post-quantum cryptography standardization,'' NIST, 2024. [Online]. Available: \url{https://csrc.nist.gov/projects/post-quantum-cryptography}

\bibitem{kyber2023}
R.~Avanzi \textit{et~al.}, ``CRYSTALS-Kyber: A CCA-secure module-lattice-based KEM,'' in \textit{Proc. IEEE Symp. Secur. Privacy}, 2023. [Online]. Available: \url{https://doi.org/10.1109/SP46215.2023.10159101}

\bibitem{dilithium2022}
S.~Bai \textit{et~al.}, ``CRYSTALS-Dilithium: A lattice-based digital signature scheme,'' \textit{IACR Trans. Cryptogr. Hardw. Embed. Syst.}, vol.~2022, no.~1, pp.~218--244, 2022. [Online]. Available: \url{https://doi.org/10.46586/tches.v2022.i1.218-244}

\bibitem{bernstein2010}
D.~J. Bernstein, ``Grover vs.\ McEliece,'' in \textit{Post-Quantum Cryptography}, LNCS, vol.~6061. Springer, 2010, pp.~89--98. [Online]. Available: \url{https://doi.org/10.1007/978-3-642-12929-2_7}

\bibitem{alagic2022}
G.~Alagic \textit{et~al.}, ``Status report on the third round of the NIST post-quantum cryptography standardization process,'' NIST IR 8413, NIST, Gaithersburg, MD, USA, Sep. 2022. [Online]. Available: \url{https://doi.org/10.6028/NIST.IR.8413-upd1}

\bibitem{menezes2001}
A.~J. Menezes, P.~C.~van Oorschot, and S.~A. Vanstone, \textit{Handbook of Applied Cryptography}, 5th~ed.\ Boca Raton, FL, USA: CRC Press, 2001. [Online]. Available: \url{https://cacr.uwaterloo.ca/hac/}

\bibitem{shoup2008}
V.~Shoup, \textit{A Computational Introduction to Number Theory and Algebra}, Version~2.3.1. Cambridge, U.K.: Cambridge Univ.\ Press, 2008. [Online]. Available: \url{https://shoup.net/ntb/}

\bibitem{kelsey1999}
J.~A. Kelsey, B.~Schneier, and N.~Ferguson, ``Yarrow-160: Notes on the design and analysis of the Yarrow cryptographic pseudorandom number generator,'' in \textit{Proc. Sixth Annu. Workshop Sel. Areas Cryptogr.}, 1999. [Online]. Available: \url{https://www.schneier.com/wp-content/uploads/1999/12/yarrow.pdf}

\bibitem{cryptography_io}
Cryptography Software Foundation, ``cryptography---A Python library exposing cryptographic recipes and primitives,'' Version~46.0.7, 2025. [Online]. Available: \url{https://cryptography.io/}

\bibitem{python_tracemalloc}
Python Software Foundation, ``The Python standard library---tracemalloc: Trace memory allocations,'' Python~3 Documentation, 2025. [Online]. Available: \url{https://docs.python.org/3/library/tracemalloc.html}

\bibitem{python_threading}
Python Software Foundation, ``The Python standard library---threading: Thread-based parallelism,'' Python~3 Documentation, 2025. [Online]. Available: \url{https://docs.python.org/3/library/threading.html}

\end{thebibliography}


\begin{IEEEbiographynophoto}{Bhavesh Patil}
is an undergraduate student in the Department of Information Technology at C.~K.~Pithawala College of Engineering, Surat, Gujarat, India. His research interests include cryptography, network security, and web application performance.
\end{IEEEbiographynophoto}

\begin{IEEEbiographynophoto}{Devendra Patil}
is an undergraduate student in the Department of Information Technology at C.~K.~Pithawala College of Engineering, Surat, Gujarat, India. His research interests include cybersecurity and cryptographic algorithms.
\end{IEEEbiographynophoto}

\begin{IEEEbiographynophoto}{Sakti Vala}
is an undergraduate student in the Department of Information Technology at C.~K.~Pithawala College of Engineering, Surat, Gujarat, India. His research interests include information security and distributed systems.
\end{IEEEbiographynophoto}

\end{document}
